% DSE demo review, freeze candidate. Numbers, checklist results, the scope list, the demo script and the register rows
% are generated by build_review.py into generated/*.tex from the dashboard API, the register, the ledgers and git;
% nothing numeric is typed here. Build (in demo/review/):
%   py -3.8 capture.py shots ; py -3.12 build_review.py --tests "<pytest summary>" ; pdflatex demo_review.tex (twice)
\documentclass[11pt]{article}
\usepackage[T1]{fontenc}
\usepackage{lmodern}
\usepackage[a4paper,margin=2.2cm]{geometry}
\usepackage{graphicx}
\usepackage{booktabs}
\usepackage{longtable}
\usepackage{array}
\usepackage[font=small,labelfont=bf]{caption}
\usepackage{float}
\usepackage{enumitem}
\usepackage{microtype}
\setlist[itemize]{topsep=2pt,itemsep=1pt}
\setlength{\LTcapwidth}{\textwidth}
\newcommand{\metaBranch}{demo/oct28}
\newcommand{\metaCommit}{73e9fc7}
\newcommand{\metaDate}{2026-10-22}
\newcommand{\metaTests}{not run}
\newcommand{\metaFails}{1}
\newcommand{\metaNchecks}{24}

\newcommand{\shot}[3]{\begin{figure}[H]\centering
  \includegraphics[width=\textwidth,height=0.62\textheight,keepaspectratio]{shots/#1.png}
  \caption{#2}\label{#3}\end{figure}}

\begin{document}

\begin{titlepage}
\centering
\vspace*{5cm}
{\Huge\bfseries DSE Demo Review\par}
\vspace{0.4cm}
{\LARGE freeze candidate\par}
\vspace{1.5cm}
{\large Rithwik Sharma\par}
\vspace{0.5cm}
{\large \metaDate\par}
\vspace{1.5cm}
\begin{tabular}{rl}
Branch & \texttt{\metaBranch}\\
Commit & \texttt{\metaCommit} (content reviewed; this PDF is committed on top of it)\\
Demo freeze & 2026-10-28\\
Test suite & \metaTests\\
Checklist & \metaNchecks{} checks, \metaFails{} failed\\
\end{tabular}
\vfill
{\small Every number in this document is generated from the dashboard's offline cache (built from
docs/NUMBERS\_REGISTER.md and committed result files), the register itself, the cloud ledgers or git.\par}
\end{titlepage}

\tableofcontents
\clearpage

\section{Summary}
What the demo shows:
\begin{itemize}
\item \textbf{The failure is silent.} On the same evo-t2s machine and the same models, the Intel CPU device makes
  Ollama load a 4096-token window and every session loses a rule after the window fills, while the iGPU loads 32768;
  no call surfaces an error. On evo-x2 the qwen3:14b wrong-task rate rises from the default window to 4096 with every
  wrong result returned as a normal answer (view A, Table~\ref{tab:a}).
\item \textbf{The obvious fixes do not hold.} num\_keep reduces but does not remove the wrong-task excess; client-side
  trimming breaks a rule the untrimmed run kept; a bigger memory setting is up to 16x slower or does not start (story
  step 3).
\item \textbf{The tool decides at session start, per device and model}: route local only if local accuracy is at least
  all-cloud accuracy minus a stated tolerance, else send the whole session to the cloud from the start; it shows the
  recommended local setting (window, num\_keep, KV precision, flash attention, each labelled) and the tolerance at
  which each device would route local. On evo-x2 the KV q8\_0 / q4\_0 settings are measured and attributed to KV precision by the verified
  flash-attention control; on evo-t2s ``kept local'' is measured by R2 session pass rate, not B3 task accuracy
  (view B, Table~\ref{tab:b}).
\item \textbf{The money:} story step 4 shows USD 2.53 per long task in the cloud (priced) and about 2 cents to fix a
  wrong answer later (measured; fix rate, 95\% Wilson CI and the inherited-actions caveat in its Details), then the
  accuracy-tolerance slider: how much accuracy gap you accept decides where it runs. Every dollar figure that depends
  on volume says ``at 1,000 tasks a day'' and is tagged projected (measured cost per task times an assumed volume).
  The ``Fix one wrong answer now (live)'' button sends the view D task through the capped client with the
  full-history protocol (one live call per demo run) and falls back to that task's recorded redo, tagged recorded
  (Figure~\ref{fig:s4l}). The measured hybrid trade-off,
  with the ceiling only as a side note, is in its Details and in view C (Table~\ref{tab:c}).
\item \textbf{One live escalation} (engineering view D) of a real wrong task through the capped client, with the
  task's recorded full-history redo as the fallback (Table~\ref{tab:d}).
\end{itemize}
What it deliberately does not claim:
\begin{itemize}
\item \textbf{No router savings.} The register row b5-verdict is refuted on evo-x2 (the router policy was dominated), so
  no savings figure is shown anywhere, including the older recommendation screen.
\item \textbf{No figure from the first cloud redo.} results/cloud\_redo\_v1.jsonl was a broken protocol run (thinking
  consumed the output budget); its row and the impact rows built on it are withheld, and a test checks that none of
  its figures appears. The full-history redo ran with up to 2048 output tokens per call against the local run's 384,
  and the screen says so.
\item \textbf{No detector ceiling as a headline.} The perfect-detector figure appears only as a dotted line and a side
  note labelled ``ceiling, no such detector''.
\item \textbf{No local route at matched quality.} The slider and the tool card follow row hybrid-route-at-start: at
  tolerance 0 every measured device routes the whole session to the cloud from the start; at the default of 5 points
  too; the AMD workstation routes local from 7.8 (Haiku) / 11.1 (Sonnet) points, the Intel laptop from 16.7 / 20.0
  (R2 session pass rate, not B3 task accuracy).
\item \textbf{Story wording.} Step 1's green card says ``kept their rules'', not ``ran correctly'': those sessions lost
  no gated rule, but their history recall did not survive every checkpoint.
\end{itemize}

\section{Story mode}
The dashboard lands in story mode, written for a non-expert audience: five steps, one idea and one big number each,
plain words (Intel laptop, AMD workstation, NVIDIA laptop; memory setting; wrong answers), presenter-size type, one
small tag per number, and every source in a closed ``Details'' panel per step (Figures~\ref{fig:s1} to
\ref{fig:s5b}). The engineering view (Section~\ref{sec:eng} and the view sections) is behind a single link, top
right. Cards and type are sized for a 1600x1000 full-screen view. The story runs entirely from the cache, apart from
the one optional live call in step 4. A three-minute backup recording (no audio, no live call; the step 4 button shows
its recorded fallback) is
\texttt{demo/review/demo\_story.webm}; the two-page leave-behind is \texttt{demo/review/LEAVE\_BEHIND.pdf}; the
rehearsal sheet is \texttt{demo/review/STEP\_SCRIPT.md}.
\shot{story_step1}{Step 1, the landing page: same laptop, two settings.}{fig:s1}
\shot{story_step2}{Step 2: how often answers are silently wrong (rounded on screen; exact counts in Details).}{fig:s2}
\shot{story_step2_details}{Step 2 with its Details panel opened: the exact counts and register row.}{fig:s2d}
\shot{story_step3}{Step 3: why not just change one setting.}{fig:s3}
\shot{story_step4}{Step 4: what it costs, with the accuracy-tolerance slider at its default of 5 points.}{fig:s4}
\shot{story_step4_t12}{Step 4 with the slider at 12 points: the AMD workstation routes local.}{fig:s4b}
\shot{story_step4_live}{Step 4 after ``Fix one wrong answer now (live)'' with no live connection: the recorded
full-history redo of the same task, tagged recorded.}{fig:s4l}
\shot{story_step5}{Step 5: what our tool does, on the AMD workstation at the default gap.}{fig:s5}
\shot{story_step5_intel}{Step 5 on the Intel laptop (local accuracy measured by R2 session pass rate).}{fig:s5b}

\section{Engineering view}\label{sec:eng}
Behind the ``Engineering view'' link: the four views below and the frontier and replay screens. Cloud rows use the
Anthropic shared price table and the measured cloud rows (no OpenAI stub rows); hardware with no measured envelope is
listed in one ``Not measured'' line; every verbatim register text sits in a closed ``Details'' disclosure under a
one-sentence summary.
\shot{eng_top}{Engineering view, top: banner with the Anthropic price table, then view A.}{fig:e1}
\shot{eng_inputs}{Engineering view: inputs, with the one-line ``Not measured'' hardware list.}{fig:e1b}
\shot{eng_frontier}{Engineering view: frontier per machine and the measured Anthropic cloud rows.}{fig:e2}
\shot{eng_replay}{Engineering view: the replay screen, register rows summarised above closed Details.}{fig:e3}

\section{View A: silent failure}
\textbf{Takeaway.} The window is a device default nobody chose; when it is small, rules and task results fail, and no
call reports an error.
\shot{A_default}{View A, default state.}{fig:a1}
\shot{A_details}{View A after opening the per-model, per-window rule-loss grid (strengthened run v1 + v1b).}{fig:a2}
\begin{longtable}{>{\raggedright\arraybackslash}p{0.54\textwidth}>{\raggedright\arraybackslash}p{0.25\textwidth}>{\raggedright\arraybackslash}p{0.13\textwidth}}
\caption{View A: every number shown}\label{tab:a}\\
\toprule
\textbf{Value shown} & \textbf{Register id / file} & \textbf{Status} \\
\midrule
\endfirsthead
\toprule
\textbf{Value shown} & \textbf{Register id / file} & \textbf{Status} \\
\midrule
\endhead
\small evo-t2s Intel CPU, llama3.1:8b: loaded window [4096] & {\footnotesize \texttt{t2s-\allowbreak{}r2-\allowbreak{}real-\allowbreak{}cpu}} & \small measured \\

\small evo-t2s Intel CPU, llama3.1:8b: 5/5 sessions silently lost a rule, 5 after the window filled & {\footnotesize \texttt{t2s-\allowbreak{}r2-\allowbreak{}real-\allowbreak{}cpu}} & \small measured \\

\small evo-t2s Intel CPU, qwen3:14b: loaded window [4096] & {\footnotesize \texttt{t2s-\allowbreak{}r2-\allowbreak{}real-\allowbreak{}cpu}} & \small measured \\

\small evo-t2s Intel CPU, qwen3:14b: 5/5 sessions silently lost a rule, 5 after the window filled & {\footnotesize \texttt{t2s-\allowbreak{}r2-\allowbreak{}real-\allowbreak{}cpu}} & \small measured \\

\small evo-t2s Intel CPU: 0 of 800 calls surfaced an error & {\footnotesize \texttt{results/\allowbreak{}t2s\_\allowbreak{}r2\_\allowbreak{}real\_\allowbreak{}cpu\_\allowbreak{}v1.\allowbreak{}jsonl}} & \small measured \\

\small evo-t2s Intel iGPU, llama3.1:8b: loaded window [32768] & {\footnotesize \texttt{t2s-\allowbreak{}r2-\allowbreak{}real-\allowbreak{}igpu}} & \small measured \\

\small evo-t2s Intel iGPU, llama3.1:8b: 0/5 sessions silently lost a rule, 0 after the window filled & {\footnotesize \texttt{t2s-\allowbreak{}r2-\allowbreak{}real-\allowbreak{}igpu}} & \small measured \\

\small evo-t2s Intel iGPU, qwen3:14b: loaded window [32768] & {\footnotesize \texttt{t2s-\allowbreak{}r2-\allowbreak{}real-\allowbreak{}igpu}} & \small measured \\

\small evo-t2s Intel iGPU, qwen3:14b: 1/5 sessions silently lost a rule, 0 after the window filled & {\footnotesize \texttt{t2s-\allowbreak{}r2-\allowbreak{}real-\allowbreak{}igpu}} & \small measured \\

\small evo-t2s Intel iGPU: 0 of 800 calls surfaced an error & {\footnotesize \texttt{results/\allowbreak{}t2s\_\allowbreak{}r2\_\allowbreak{}real\_\allowbreak{}igpu\_\allowbreak{}v1.\allowbreak{}jsonl}} & \small measured \\

\small qwen3:14b ollama\_default: wrong 17/150 (11.3\%), silent 17, surfaced 0 & {\footnotesize \texttt{x2-\allowbreak{}evalset-\allowbreak{}tasks-\allowbreak{}real-\allowbreak{}v2-\allowbreak{}wrong-\allowbreak{}rate}} & \small measured \\

\small qwen3:14b ollama\_ctx\_16384: wrong 19/150 (12.7\%), silent 19, surfaced 0 & {\footnotesize \texttt{x2-\allowbreak{}evalset-\allowbreak{}tasks-\allowbreak{}real-\allowbreak{}v2-\allowbreak{}wrong-\allowbreak{}rate}} & \small measured \\

\small ollama\_ctx\_16384: excess +1.3 points (95\% CI -1.3 to +4.0) & {\footnotesize \texttt{x2-\allowbreak{}evalset-\allowbreak{}tasks-\allowbreak{}real-\allowbreak{}v2-\allowbreak{}wrong-\allowbreak{}rate}} & \small measured \\

\small qwen3:14b ollama\_ctx\_8192: wrong 53/150 (35.3\%), silent 53, surfaced 0 & {\footnotesize \texttt{x2-\allowbreak{}evalset-\allowbreak{}tasks-\allowbreak{}real-\allowbreak{}v2-\allowbreak{}wrong-\allowbreak{}rate}} & \small measured \\

\small ollama\_ctx\_8192: excess +24.0 points (95\% CI +18.7 to +29.3) & {\footnotesize \texttt{x2-\allowbreak{}evalset-\allowbreak{}tasks-\allowbreak{}real-\allowbreak{}v2-\allowbreak{}wrong-\allowbreak{}rate}} & \small measured \\

\small qwen3:14b ollama\_ctx\_4096: wrong 63/150 (42.0\%), silent 63, surfaced 0 & {\footnotesize \texttt{x2-\allowbreak{}evalset-\allowbreak{}tasks-\allowbreak{}real-\allowbreak{}v2-\allowbreak{}wrong-\allowbreak{}rate}} & \small measured \\

\small ollama\_ctx\_4096: excess +30.7 points (95\% CI +26.0 to +34.7) & {\footnotesize \texttt{x2-\allowbreak{}evalset-\allowbreak{}tasks-\allowbreak{}real-\allowbreak{}v2-\allowbreak{}wrong-\allowbreak{}rate}} & \small measured \\

\small qwen3:14b ollama\_ctx\_4096\_numkeep: wrong 45/150 (30.0\%), silent 45, surfaced 0 & {\footnotesize \texttt{x2-\allowbreak{}evalset-\allowbreak{}tasks-\allowbreak{}real-\allowbreak{}v2-\allowbreak{}wrong-\allowbreak{}rate}} & \small measured \\

\small ollama\_ctx\_4096\_numkeep: excess +18.7 points (95\% CI +16.0 to +21.3) & {\footnotesize \texttt{x2-\allowbreak{}evalset-\allowbreak{}tasks-\allowbreak{}real-\allowbreak{}v2-\allowbreak{}wrong-\allowbreak{}rate}} & \small measured \\

\small evo-x2 ollama\_default llama3.1:8b: 3/5 silent, 0 after window & {\footnotesize \texttt{R2-\allowbreak{}v1b-\allowbreak{}gated-\allowbreak{}per-\allowbreak{}model}} & \small measured \\

\small evo-x2 ollama\_default qwen3:14b: 0/5 silent, 0 after window & {\footnotesize \texttt{R2-\allowbreak{}v1b-\allowbreak{}gated-\allowbreak{}per-\allowbreak{}model}} & \small measured \\

\small evo-x2 ollama\_default qwen3:8b: 0/5 silent, 0 after window & {\footnotesize \texttt{R2-\allowbreak{}v1b-\allowbreak{}gated-\allowbreak{}per-\allowbreak{}model}} & \small measured \\

\small evo-x2 num\_ctx\_32768 llama3.1:8b: 1/5 silent, 0 after window & {\footnotesize \texttt{R2-\allowbreak{}v1b-\allowbreak{}gated-\allowbreak{}per-\allowbreak{}model}} & \small measured \\

\small evo-x2 num\_ctx\_32768 qwen3:14b: 0/5 silent, 0 after window & {\footnotesize \texttt{R2-\allowbreak{}v1b-\allowbreak{}gated-\allowbreak{}per-\allowbreak{}model}} & \small measured \\

\small evo-x2 num\_ctx\_32768 qwen3:8b: 0/5 silent, 0 after window & {\footnotesize \texttt{R2-\allowbreak{}v1b-\allowbreak{}gated-\allowbreak{}per-\allowbreak{}model}} & \small measured \\

\small evo-x2 num\_ctx\_16384 llama3.1:8b: 2/5 silent, 2 after window & {\footnotesize \texttt{R2-\allowbreak{}v1b-\allowbreak{}gated-\allowbreak{}per-\allowbreak{}model}} & \small measured \\

\small evo-x2 num\_ctx\_16384 qwen3:14b: 0/5 silent, 0 after window & {\footnotesize \texttt{R2-\allowbreak{}v1b-\allowbreak{}gated-\allowbreak{}per-\allowbreak{}model}} & \small measured \\

\small evo-x2 num\_ctx\_16384 qwen3:8b: 0/5 silent, 0 after window & {\footnotesize \texttt{R2-\allowbreak{}v1b-\allowbreak{}gated-\allowbreak{}per-\allowbreak{}model}} & \small measured \\

\small evo-x2 num\_ctx\_8192 llama3.1:8b: 4/5 silent, 4 after window & {\footnotesize \texttt{R2-\allowbreak{}v1b-\allowbreak{}gated-\allowbreak{}per-\allowbreak{}model}} & \small measured \\

\small evo-x2 num\_ctx\_8192 qwen3:14b: 0/5 silent, 0 after window & {\footnotesize \texttt{R2-\allowbreak{}v1b-\allowbreak{}gated-\allowbreak{}per-\allowbreak{}model}} & \small measured \\

\small evo-x2 num\_ctx\_8192 qwen3:8b: 4/5 silent, 4 after window & {\footnotesize \texttt{R2-\allowbreak{}v1b-\allowbreak{}gated-\allowbreak{}per-\allowbreak{}model}} & \small measured \\

\small evo-x2 num\_ctx\_4096 llama3.1:8b: 5/5 silent, 5 after window & {\footnotesize \texttt{R2-\allowbreak{}v1b-\allowbreak{}gated-\allowbreak{}per-\allowbreak{}model}} & \small measured \\

\small evo-x2 num\_ctx\_4096 qwen3:14b: 4/5 silent, 4 after window & {\footnotesize \texttt{R2-\allowbreak{}v1b-\allowbreak{}gated-\allowbreak{}per-\allowbreak{}model}} & \small measured \\

\small evo-x2 num\_ctx\_4096 qwen3:8b: 1/5 silent, 1 after window & {\footnotesize \texttt{R2-\allowbreak{}v1b-\allowbreak{}gated-\allowbreak{}per-\allowbreak{}model}} & \small measured \\

\bottomrule
\end{longtable}


\section{View B: per-device recommendation}
\textbf{Takeaway.} The decision at session start, and the recommended window, num\_keep, KV precision and flash
attention with their measured outcome, or ``not measured''.
\shot{B_default}{View B, default state: evo-x2 / qwen3:14b / multi-turn agent session.}{fig:b1}
\shot{B_cpu}{View B with the device switched to the evo-t2s Intel CPU: no measured setting holds the rules for
qwen3:14b, so the whole session goes to the cloud; kept local is measured by R2 session pass rate.}{fig:b2}
\shot{B_blade}{View B with the device switched to the Blade: every field not measured.}{fig:b3}
\begin{longtable}{>{\raggedright\arraybackslash}p{0.54\textwidth}>{\raggedright\arraybackslash}p{0.25\textwidth}>{\raggedright\arraybackslash}p{0.13\textwidth}}
\caption{View B: every number shown, for the eight device / model / workload choices reviewed}\label{tab:b}\\
\toprule
\textbf{Value shown} & \textbf{Register id / file} & \textbf{Status} \\
\midrule
\endfirsthead
\toprule
\textbf{Value shown} & \textbf{Register id / file} & \textbf{Status} \\
\midrule
\endhead
\small evo-x2 / qwen3:14b / r2\_agent: window = 40960 & {\footnotesize \texttt{results/\allowbreak{}x2\_\allowbreak{}policy\_\allowbreak{}compare.\allowbreak{}jsonl:\allowbreak{}8}} & \small measured \\

\small evo-x2 / qwen3:14b / r2\_agent: kv\_precision = q8\_0 & {\footnotesize \texttt{results/\allowbreak{}x2\_\allowbreak{}policy\_\allowbreak{}compare.\allowbreak{}jsonl:\allowbreak{}4}} & \small measured \\

\small evo-x2 / qwen3:14b / r2\_agent: flash\_attention = on (OLLAMA\_FLASH\_ATTENTION=1) & {\footnotesize \texttt{results/\allowbreak{}x2\_\allowbreak{}policy\_\allowbreak{}compare.\allowbreak{}jsonl:\allowbreak{}4}} & \small measured \\

\small evo-x2 / qwen3:14b / r2\_agent: num\_keep = not set & {\footnotesize \texttt{results/\allowbreak{}x2\_\allowbreak{}policy\_\allowbreak{}compare.\allowbreak{}jsonl:\allowbreak{}8}} & \small measured \\

\small evo-x2 / qwen3:14b / r2\_agent: context\_handling = full history & {\footnotesize \texttt{results/\allowbreak{}x2\_\allowbreak{}policy\_\allowbreak{}compare.\allowbreak{}jsonl:\allowbreak{}8}} & \small measured \\

\small evo-x2 / qwen3:14b / r2\_agent: rules held 100.0\%, history recalled 62.5\% & {\footnotesize \texttt{b5-\allowbreak{}policy-\allowbreak{}metrics}} & \small measured \\

\small evo-x2 / qwen3:14b / r2\_agent: 75.4 min per 40-turn session & {\footnotesize \texttt{b5-\allowbreak{}policy-\allowbreak{}metrics}} & \small measured \\

\small evo-x2 / qwen3:14b / r2\_agent: default (a): 214.8 min, rules 100.0\%, history 62.5\%; saves 139.4 min per session & {\footnotesize \texttt{b5-\allowbreak{}policy-\allowbreak{}metrics}} & \small measured \\

\small evo-x2 / qwen3:14b / r2\_agent: long-context recall probe at this precision: 100.0\% at ctx 40960 & {\footnotesize \texttt{b5-\allowbreak{}kv-\allowbreak{}recall}} & \small measured \\

\small evo-x2 / qwen3:14b / r2\_agent: Decide at session start: at matched quality (tolerance 0) the measured answer is the whole session to the cloud from the start. & {\footnotesize \texttt{hybrid-\allowbreak{}route-\allowbreak{}at-\allowbreak{}start}} & \small measured \\

\small evo-x2 / qwen3:14b / r2\_agent: would route local only at tolerance \textgreater{}= 7.8 points vs claude-haiku-5-5 (local 80/90 vs all-cloud 87/90; measured, B3 task accuracy) & {\footnotesize \texttt{hybrid-\allowbreak{}route-\allowbreak{}at-\allowbreak{}start}} & \small measured \\

\small evo-x2 / qwen3:14b / r2\_agent: would route local only at tolerance \textgreater{}= 11.1 points vs claude-sonnet-5-5 (local 80/90 vs all-cloud 90/90; measured, B3 task accuracy) & {\footnotesize \texttt{hybrid-\allowbreak{}route-\allowbreak{}at-\allowbreak{}start}} & \small measured \\

\small evo-x2 / llama3.1:8b / r2\_agent: window = 131072 & {\footnotesize \texttt{results/\allowbreak{}x2\_\allowbreak{}policy\_\allowbreak{}compare.\allowbreak{}jsonl:\allowbreak{}7}} & \small measured \\

\small evo-x2 / llama3.1:8b / r2\_agent: kv\_precision = q4\_0 & {\footnotesize \texttt{results/\allowbreak{}x2\_\allowbreak{}policy\_\allowbreak{}compare.\allowbreak{}jsonl:\allowbreak{}5}} & \small measured \\

\small evo-x2 / llama3.1:8b / r2\_agent: flash\_attention = on (OLLAMA\_FLASH\_ATTENTION=1) & {\footnotesize \texttt{results/\allowbreak{}x2\_\allowbreak{}policy\_\allowbreak{}compare.\allowbreak{}jsonl:\allowbreak{}5}} & \small measured \\

\small evo-x2 / llama3.1:8b / r2\_agent: num\_keep = not set & {\footnotesize \texttt{results/\allowbreak{}x2\_\allowbreak{}policy\_\allowbreak{}compare.\allowbreak{}jsonl:\allowbreak{}7}} & \small measured \\

\small evo-x2 / llama3.1:8b / r2\_agent: context\_handling = full history & {\footnotesize \texttt{results/\allowbreak{}x2\_\allowbreak{}policy\_\allowbreak{}compare.\allowbreak{}jsonl:\allowbreak{}7}} & \small measured \\

\small evo-x2 / llama3.1:8b / r2\_agent: rules held 100.0\%, history recalled 100.0\% & {\footnotesize \texttt{b5-\allowbreak{}policy-\allowbreak{}metrics}} & \small measured \\

\small evo-x2 / llama3.1:8b / r2\_agent: 6.6 min per 40-turn session & {\footnotesize \texttt{b5-\allowbreak{}policy-\allowbreak{}metrics}} & \small measured \\

\small evo-x2 / llama3.1:8b / r2\_agent: default (a): 43.7 min, rules 100.0\%, history 95.8\%; saves 37.1 min per session & {\footnotesize \texttt{b5-\allowbreak{}policy-\allowbreak{}metrics}} & \small measured \\

\small evo-x2 / llama3.1:8b / r2\_agent: long-context recall probe at this precision: 58.3\% at ctx 131072 & {\footnotesize \texttt{b5-\allowbreak{}kv-\allowbreak{}recall}} & \small measured \\

\small evo-x2 / llama3.1:8b / r2\_agent: Decide at session start: run locally with the recommended setting. & {\footnotesize \texttt{b5-\allowbreak{}policy-\allowbreak{}metrics, b5-\allowbreak{}kv-\allowbreak{}recall, b5-\allowbreak{}fa-\allowbreak{}control, cost-\allowbreak{}cut-\allowbreak{}config-\allowbreak{}time}} & \small measured \\

\small evo-x2 / llama3.1:8b / r2\_agent: route at start: not measured for this device and model (row hybrid-route-at-start covers qwen3:14b on evo-x2 and the evo-t2s iGPU) & {\footnotesize \texttt{none}} & \small not measured \\

\small evo-t2s-cpu / qwen3:14b / r2\_agent: window = not measured & {\footnotesize \texttt{none}} & \small not measured \\

\small evo-t2s-cpu / qwen3:14b / r2\_agent: Ollama default window 4096: 5/5 sessions silently lost a gated rule, 5 after the window filled & {\footnotesize \texttt{t2s-\allowbreak{}r2-\allowbreak{}real-\allowbreak{}cpu}} & \small measured \\

\small evo-t2s-cpu / qwen3:14b / r2\_agent: with client-side trimming: not recovered (pre-registered rule) & {\footnotesize \texttt{t2s-\allowbreak{}r2-\allowbreak{}mitigation-\allowbreak{}cpu}} & \small measured \\

\small evo-t2s-cpu / qwen3:14b / r2\_agent: kept local, measured by R2 session pass rate (not B3 task accuracy): 0/5 sessions & {\footnotesize \texttt{t2s-\allowbreak{}config-\allowbreak{}comparison}} & \small measured \\

\small evo-t2s-cpu / qwen3:14b / r2\_agent: saves USD 0.00/month vs claude-haiku-5-5 at 1,000 tasks/day (kept local share x bill) & {\footnotesize \texttt{t2s-\allowbreak{}config-\allowbreak{}comparison}} & \small projected \\

\small evo-t2s-cpu / qwen3:14b / r2\_agent: saves USD -374.93/month vs claude-sonnet-5-5 at 1,000 tasks/day (costs more than all-cloud) & {\footnotesize \texttt{t2s-\allowbreak{}config-\allowbreak{}comparison}} & \small projected \\

\small evo-t2s-cpu / qwen3:14b / r2\_agent: overall status & {\footnotesize \texttt{t2s-\allowbreak{}r2-\allowbreak{}real-\allowbreak{}cpu, t2s-\allowbreak{}r2-\allowbreak{}mitigation-\allowbreak{}cpu, t2s-\allowbreak{}config-\allowbreak{}comparison}} & \small not measured \\

\small evo-t2s-cpu / qwen3:14b / r2\_agent: Decide at session start: send the whole session to the cloud from the start. & {\footnotesize \texttt{t2s-\allowbreak{}r2-\allowbreak{}real-\allowbreak{}cpu, t2s-\allowbreak{}r2-\allowbreak{}mitigation-\allowbreak{}cpu, t2s-\allowbreak{}config-\allowbreak{}comparison}} & \small measured \\

\small evo-t2s-cpu / qwen3:14b / r2\_agent: route at start: not measured for this device and model (row hybrid-route-at-start covers qwen3:14b on evo-x2 and the evo-t2s iGPU) & {\footnotesize \texttt{none}} & \small not measured \\

\small evo-t2s-cpu / llama3.1:8b / r2\_agent: window = 4096 & {\footnotesize \texttt{t2s-\allowbreak{}r2-\allowbreak{}real-\allowbreak{}cpu}} & \small measured \\

\small evo-t2s-cpu / llama3.1:8b / r2\_agent: context\_handling = client-side trimming that keeps the system prompt & {\footnotesize \texttt{t2s-\allowbreak{}r2-\allowbreak{}mitigation-\allowbreak{}cpu}} & \small measured \\

\small evo-t2s-cpu / llama3.1:8b / r2\_agent: Ollama default window 4096: 5/5 sessions silently lost a gated rule, 5 after the window filled & {\footnotesize \texttt{t2s-\allowbreak{}r2-\allowbreak{}real-\allowbreak{}cpu}} & \small measured \\

\small evo-t2s-cpu / llama3.1:8b / r2\_agent: with client-side trimming: recovered (pre-registered rule) & {\footnotesize \texttt{t2s-\allowbreak{}r2-\allowbreak{}mitigation-\allowbreak{}cpu}} & \small measured \\

\small evo-t2s-cpu / llama3.1:8b / r2\_agent: kept local, measured by R2 session pass rate (not B3 task accuracy): 0/5 sessions & {\footnotesize \texttt{t2s-\allowbreak{}config-\allowbreak{}comparison}} & \small measured \\

\small evo-t2s-cpu / llama3.1:8b / r2\_agent: saves USD 0.00/month vs claude-haiku-5-5 at 1,000 tasks/day (kept local share x bill) & {\footnotesize \texttt{t2s-\allowbreak{}config-\allowbreak{}comparison}} & \small projected \\

\small evo-t2s-cpu / llama3.1:8b / r2\_agent: saves USD -374.93/month vs claude-sonnet-5-5 at 1,000 tasks/day (costs more than all-cloud) & {\footnotesize \texttt{t2s-\allowbreak{}config-\allowbreak{}comparison}} & \small projected \\

\small evo-t2s-cpu / llama3.1:8b / r2\_agent: Decide at session start: run locally with the recommended setting. & {\footnotesize \texttt{t2s-\allowbreak{}r2-\allowbreak{}real-\allowbreak{}cpu, t2s-\allowbreak{}r2-\allowbreak{}mitigation-\allowbreak{}cpu, t2s-\allowbreak{}config-\allowbreak{}comparison}} & \small measured \\

\small evo-t2s-cpu / llama3.1:8b / r2\_agent: route at start: not measured for this device and model (row hybrid-route-at-start covers qwen3:14b on evo-x2 and the evo-t2s iGPU) & {\footnotesize \texttt{none}} & \small not measured \\

\small evo-t2s-igpu / qwen3:14b / r2\_agent: window = 32768 & {\footnotesize \texttt{t2s-\allowbreak{}r2-\allowbreak{}real-\allowbreak{}igpu}} & \small measured \\

\small evo-t2s-igpu / qwen3:14b / r2\_agent: context\_handling = full history & {\footnotesize \texttt{t2s-\allowbreak{}r2-\allowbreak{}real-\allowbreak{}igpu}} & \small measured \\

\small evo-t2s-igpu / qwen3:14b / r2\_agent: Ollama default window 32768: 1/5 sessions silently lost a gated rule, 0 after the window filled & {\footnotesize \texttt{t2s-\allowbreak{}r2-\allowbreak{}real-\allowbreak{}igpu}} & \small measured \\

\small evo-t2s-igpu / qwen3:14b / r2\_agent: kept local, measured by R2 session pass rate (not B3 task accuracy): 4/5 sessions & {\footnotesize \texttt{t2s-\allowbreak{}config-\allowbreak{}comparison}} & \small measured \\

\small evo-t2s-igpu / qwen3:14b / r2\_agent: saves USD 16.10/month vs claude-haiku-5-5 at 1,000 tasks/day (kept local share x bill) & {\footnotesize \texttt{t2s-\allowbreak{}config-\allowbreak{}comparison}} & \small projected \\

\small evo-t2s-igpu / qwen3:14b / r2\_agent: saves USD 169.89/month vs claude-sonnet-5-5 at 1,000 tasks/day (kept local share x bill) & {\footnotesize \texttt{t2s-\allowbreak{}config-\allowbreak{}comparison}} & \small projected \\

\small evo-t2s-igpu / qwen3:14b / r2\_agent: Decide at session start: at matched quality (tolerance 0) the measured answer is the whole session to the cloud from the start. & {\footnotesize \texttt{hybrid-\allowbreak{}route-\allowbreak{}at-\allowbreak{}start}} & \small measured \\

\small evo-t2s-igpu / qwen3:14b / r2\_agent: would route local only at tolerance \textgreater{}= 16.7 points vs claude-haiku-5-5 (local 4/5 vs all-cloud 87/90; measured by R2 session pass rate (not B3 task accuracy)) & {\footnotesize \texttt{hybrid-\allowbreak{}route-\allowbreak{}at-\allowbreak{}start}} & \small measured \\

\small evo-t2s-igpu / qwen3:14b / r2\_agent: would route local only at tolerance \textgreater{}= 20.0 points vs claude-sonnet-5-5 (local 4/5 vs all-cloud 90/90; measured by R2 session pass rate (not B3 task accuracy)) & {\footnotesize \texttt{hybrid-\allowbreak{}route-\allowbreak{}at-\allowbreak{}start}} & \small measured \\

\small evo-x2 / qwen3:14b / b3\_tasks: window = 40960 & {\footnotesize \texttt{x2-\allowbreak{}evalset-\allowbreak{}tasks-\allowbreak{}real-\allowbreak{}v2-\allowbreak{}wrong-\allowbreak{}rate}} & \small measured \\

\small evo-x2 / qwen3:14b / b3\_tasks: num\_keep = not set & {\footnotesize \texttt{x2-\allowbreak{}evalset-\allowbreak{}tasks-\allowbreak{}real-\allowbreak{}v2-\allowbreak{}wrong-\allowbreak{}rate}} & \small measured \\

\small evo-x2 / qwen3:14b / b3\_tasks: context\_handling = full history & {\footnotesize \texttt{x2-\allowbreak{}evalset-\allowbreak{}tasks-\allowbreak{}real-\allowbreak{}v2-\allowbreak{}wrong-\allowbreak{}rate}} & \small measured \\

\small evo-x2 / qwen3:14b / b3\_tasks: wrong task results 17/150 (11.3\%), 0 surfaced an error & {\footnotesize \texttt{x2-\allowbreak{}evalset-\allowbreak{}tasks-\allowbreak{}real-\allowbreak{}v2-\allowbreak{}wrong-\allowbreak{}rate}} & \small measured \\

\small evo-x2 / qwen3:14b / b3\_tasks: num\_keep at 4096: 30.0\% wrong vs 42.0\% without it; still above the best window's 11.3\% & {\footnotesize \texttt{x2-\allowbreak{}evalset-\allowbreak{}tasks-\allowbreak{}real-\allowbreak{}v2-\allowbreak{}wrong-\allowbreak{}rate}} & \small measured \\

\small evo-x2 / qwen3:14b / b3\_tasks: Decide at session start: at matched quality (tolerance 0) the measured answer is the whole session to the cloud from the start. & {\footnotesize \texttt{hybrid-\allowbreak{}route-\allowbreak{}at-\allowbreak{}start}} & \small measured \\

\small evo-x2 / qwen3:14b / b3\_tasks: would route local only at tolerance \textgreater{}= 7.8 points vs claude-haiku-5-5 (local 80/90 vs all-cloud 87/90; measured, B3 task accuracy) & {\footnotesize \texttt{hybrid-\allowbreak{}route-\allowbreak{}at-\allowbreak{}start}} & \small measured \\

\small evo-x2 / qwen3:14b / b3\_tasks: would route local only at tolerance \textgreater{}= 11.1 points vs claude-sonnet-5-5 (local 80/90 vs all-cloud 90/90; measured, B3 task accuracy) & {\footnotesize \texttt{hybrid-\allowbreak{}route-\allowbreak{}at-\allowbreak{}start}} & \small measured \\

\small evo-x2 / llama3.1:8b / single\_shot: window = runtime default & {\footnotesize \texttt{results/\allowbreak{}x2\_\allowbreak{}outcome\_\allowbreak{}table\_\allowbreak{}v3.\allowbreak{}jsonl}} & \small measured \\

\small evo-x2 / llama3.1:8b / single\_shot: context\_handling = runtime ollama\_default & {\footnotesize \texttt{results/\allowbreak{}x2\_\allowbreak{}outcome\_\allowbreak{}table\_\allowbreak{}v3.\allowbreak{}jsonl}} & \small measured \\

\small evo-x2 / llama3.1:8b / single\_shot: quality 1.000 (n 330) & {\footnotesize \texttt{results/\allowbreak{}x2\_\allowbreak{}outcome\_\allowbreak{}table\_\allowbreak{}v3.\allowbreak{}jsonl}} & \small measured \\

\small evo-x2 / llama3.1:8b / single\_shot: p50 latency 9.6 s & {\footnotesize \texttt{results/\allowbreak{}x2\_\allowbreak{}outcome\_\allowbreak{}table\_\allowbreak{}v3.\allowbreak{}jsonl}} & \small measured \\

\small evo-x2 / llama3.1:8b / single\_shot: timeout share 0.029 (floor 0.05) & {\footnotesize \texttt{results/\allowbreak{}x2\_\allowbreak{}outcome\_\allowbreak{}table\_\allowbreak{}v3.\allowbreak{}jsonl}} & \small measured \\

\small evo-x2 / llama3.1:8b / single\_shot: per call (single-shot workload): no session to route & {\footnotesize \texttt{none}} & \small context \\

\small blade / qwen3:14b / r2\_agent: window = not measured & {\footnotesize \texttt{none}} & \small not measured \\

\small blade / qwen3:14b / r2\_agent: overall status & {\footnotesize \texttt{none}} & \small not measured \\

\small blade / qwen3:14b / r2\_agent: Decide at session start: not measured for this device and model. & {\footnotesize \texttt{none}} & \small not measured \\

\small blade / qwen3:14b / r2\_agent: route at start: not measured for this device and model (row hybrid-route-at-start covers qwen3:14b on evo-x2 and the evo-t2s iGPU) & {\footnotesize \texttt{none}} & \small not measured \\

\bottomrule
\end{longtable}


\section{View C: cost panel}
\textbf{Takeaway.} Measured parts first; the lead is the measured cost of fixing a wrong answer in the cloud, with the
fix rate only alongside its CI and caveat; no router savings.
\shot{C_default}{View C, default state.}{fig:c1}
\shot{C_source}{View C after clicking the lead number: register row cloud-redo-full-context.}{fig:c2}
\begin{longtable}{>{\raggedright\arraybackslash}p{0.54\textwidth}>{\raggedright\arraybackslash}p{0.25\textwidth}>{\raggedright\arraybackslash}p{0.13\textwidth}}
\caption{View C: every number shown}\label{tab:c}\\
\toprule
\textbf{Value shown} & \textbf{Register id / file} & \textbf{Status} \\
\midrule
\endfirsthead
\toprule
\textbf{Value shown} & \textbf{Register id / file} & \textbf{Status} \\
\midrule
\endhead
\small Fixing a wrong answer in the cloud costs about USD 0.02 (mean USD 0.0201, median USD 0.0122); the cost is that the failure is silent. & {\footnotesize \texttt{cloud-\allowbreak{}redo-\allowbreak{}cost-\allowbreak{}per-\allowbreak{}wrong-\allowbreak{}task}} & \small measured \\

\small How often the answer is wrong, and silent: qwen3:14b on evo-x2: 11.3\% wrong at ollama\_default to 42.0\% at ollama\_ctx\_4096; 0 of 197 wrong results surfaced an error & {\footnotesize \texttt{x2-\allowbreak{}evalset-\allowbreak{}tasks-\allowbreak{}real-\allowbreak{}v2-\allowbreak{}wrong-\allowbreak{}rate}} & \small measured \\

\small Time saved by the measured setting, llama3.1:8b: 37.1 min per 40-turn session (85\%): default 43.7 min vs kv\_q4\_0 6.6 min (flash attention and KV precision changed together; b5-fa-control has no reading yet) & {\footnotesize \texttt{cost-\allowbreak{}cut-\allowbreak{}config-\allowbreak{}time}} & \small measured \\

\small Time saved by the measured setting, qwen3:14b: 139.4 min per 40-turn session (65\%): default 214.8 min vs kv\_q8\_0 75.4 min (flash attention and KV precision changed together; b5-fa-control has no reading yet) & {\footnotesize \texttt{cost-\allowbreak{}cut-\allowbreak{}config-\allowbreak{}time}} & \small measured \\

\small Cloud redo of each wrong task: USD 0.0201 mean, USD 0.0122 median per wrong task on claude-sonnet-5-5 (197 tasks, USD 3.95 in all; mean 6.89 s per task) & {\footnotesize \texttt{cloud-\allowbreak{}redo-\allowbreak{}cost-\allowbreak{}per-\allowbreak{}wrong-\allowbreak{}task}} & \small measured \\

\small Re-sending the history to the cloud, per session: 1,798,884 tokens re-sent; USD 0.1799 with caching, USD 3.5978 without & {\footnotesize \texttt{impact-\allowbreak{}P3}} & \small priced, not executed \\

\small Wrong answers per day for one stated team (P1): USD 5.47 to 132.31 per day at 40 sessions/day (assumed volume; time valued 0 to the full hour value; cloud part measured) & {\footnotesize \texttt{impact-\allowbreak{}v2-\allowbreak{}P1}} & \small projected \\

\small Ranking of the problems by cost (one assumed profile): P1 \textgreater{} P5 \textgreater{} P2 \textgreater{} P3 \textgreater{} P4 & {\footnotesize \texttt{impact-\allowbreak{}ranking-\allowbreak{}v2}} & \small projected \\

\small Today's cost for the reference team, measured parts only: USD 3.60 per day (measured token counts and list prices at the assumed team volume; the register calls this measured-only) & {\footnotesize \texttt{impact-\allowbreak{}total}} & \small projected \\

\small Router savings: not shown. Row b5-verdict is refuted on evo-x2 (the router policy was dominated); no register row supports a savings figure. & {\footnotesize \texttt{b5-\allowbreak{}verdict}} & \small not shown \\

\bottomrule
\end{longtable}


\section{View D: live escalation}
\textbf{Takeaway.} One real wrong task, its local answer, and its full-history cloud redo: cost, seconds and answer.
\shot{D_default}{View D, default state: the task, the local answer and the recorded full-history redo.}{fig:d1}
\shot{D_sent}{View D after ``Send to the cloud'' with live calls off: the recorded full-history redo, labelled
measured.}{fig:d2}
\begin{longtable}{>{\raggedright\arraybackslash}p{0.54\textwidth}>{\raggedright\arraybackslash}p{0.25\textwidth}>{\raggedright\arraybackslash}p{0.13\textwidth}}
\caption{View D: every number shown}\label{tab:d}\\
\toprule
\textbf{Value shown} & \textbf{Register id / file} & \textbf{Status} \\
\midrule
\endfirsthead
\toprule
\textbf{Value shown} & \textbf{Register id / file} & \textbf{Status} \\
\midrule
\endhead
\small task cal\_default\_len, qwen3:14b at ollama\_ctx\_4096, seed 20260901, turn 4 & {\footnotesize \texttt{results/\allowbreak{}x2\_\allowbreak{}evalset\_\allowbreak{}tasks\_\allowbreak{}real\_\allowbreak{}v2.\allowbreak{}jsonl:\allowbreak{}488}} & \small measured \\

\small local tool call create\_event: \{"title": "Team retro", "date": "2026-12-22", "time": "11:00", "duration\_min": 60\} & {\footnotesize \texttt{results/\allowbreak{}x2\_\allowbreak{}evalset\_\allowbreak{}tasks\_\allowbreak{}real\_\allowbreak{}v2.\allowbreak{}jsonl:\allowbreak{}488}} & \small measured \\

\small expected: \{"title": "Team retro", "date": "2026-12-22", "time": "11:00", "duration\_min": 50\} & {\footnotesize \texttt{results/\allowbreak{}x2\_\allowbreak{}evalset\_\allowbreak{}tasks\_\allowbreak{}real\_\allowbreak{}v2.\allowbreak{}jsonl:\allowbreak{}488}} & \small measured \\

\small local end state wrong: True; HTTP error: False & {\footnotesize \texttt{results/\allowbreak{}x2\_\allowbreak{}evalset\_\allowbreak{}tasks\_\allowbreak{}real\_\allowbreak{}v2.\allowbreak{}jsonl:\allowbreak{}488}} & \small measured \\

\small recorded full-history redo of this task (claude-sonnet-5-5): USD 0.0243, 4.2 s & {\footnotesize \texttt{results/\allowbreak{}cloud\_\allowbreak{}redo\_\allowbreak{}full\_\allowbreak{}context\_\allowbreak{}v1\_\allowbreak{}part3.\allowbreak{}jsonl:\allowbreak{}39}} & \small measured \\

\small spend cap for the live call: USD 10.00; live calls per run 1 & {\footnotesize \texttt{demo/\allowbreak{}dashboard/\allowbreak{}escalate.\allowbreak{}py}} & \small configuration \\

\bottomrule
\end{longtable}


\section{Review checklist}
\begin{longtable}{>{\raggedright\arraybackslash}p{0.38\textwidth}p{0.07\textwidth}>{\raggedright\arraybackslash}p{0.47\textwidth}}
\caption{Review checklist, evaluated by demo/review/build\_review.py against the dashboard API}\label{tab:check}\\\toprule
\textbf{Check} & \textbf{Result} & \textbf{Evidence} \\
\midrule\endhead
Silent failure: 5/5 on the Intel CPU default (4096) for both models & \textbf{PASS} & llama3.1:8b 5/5 at [4096]; qwen3:14b 5/5 at [4096] \\

Silent failure: qwen3:14b wrong-task rate 11.3\% to 42.0\%, 0 surfaced & \textbf{PASS} & default 11.3\%, 4096 42.0\% \\

"pending flash-attention control" shown on q8\_0 (qwen3:14b) and q4\_0 (llama3.1:8b) & \textbf{PASS} & qwen3:14b q8\_0: pending flash-attention control; llama3.1:8b q4\_0: pending flash-attention control \\

"not measured" where there is no data (Blade, every field) & \textbf{PASS} & status not measured \\

Cost panel leads with the USD 0.02 line & \textbf{PASS} & Fixing a wrong answer in the cloud costs about USD 0.02 (mean USD 0.0201, median USD 0.0122); the cost is that the failure is silent. \\

Cost panel shows no cloud fix rate & \textbf{PASS} & no fix-rate field or text in the panel; the clause is withheld in the click-through register row \\

Cost panel shows no router savings & \textbf{PASS} & Router savings: not shown. Row b5-verdict is refuted on evo-x2 (the router policy was dominated); no register row supports a savings figure. \\

Live view shows the recorded redo when not sent live & \textbf{PASS} & mode recorded redo, USD 0.0122, label measured \\

Scope: every changed file vs main is in demo/, src/dse/ or the dashboard tests & \textbf{PASS} & 28 files changed \\

Full test suite & \textbf{PASS} & 2471 passed, 14 skipped \\

Cloud spend so far (read-only ledgers) & \textbf{PASS} & main results/cloud\_ledger.jsonl USD 21.1913 over 817 rows; demo ledger USD 0.0174 over 1 call(s) \\

\bottomrule
\end{longtable}

\noindent Files changed on this branch relative to main (\texttt{git diff -{}-name-only main...HEAD}):
{\footnotesize\raggedright \texttt{demo/DEMO\_SCRIPT.md}; \texttt{demo/dashboard/build\_cache.py}; \texttt{demo/dashboard/cache/cost.json}; \texttt{demo/dashboard/cache/device.json}; \texttt{demo/dashboard/cache/escalation.json}; \texttt{demo/dashboard/cache/manifest.json}; \texttt{demo/dashboard/cache/mitigation.json}; \texttt{demo/dashboard/cache/r2\_sessions.json}; \texttt{demo/dashboard/cache/register.json}; \texttt{demo/dashboard/cache/silent.json}; \texttt{demo/dashboard/cache/story.json}; \texttt{demo/dashboard/data.py}; \texttt{demo/dashboard/escalate.py}; \texttt{demo/dashboard/server.py}; \texttt{demo/dashboard/static/app.css}; \texttt{demo/dashboard/static/app.js}; \texttt{demo/dashboard/static/index.html}; \texttt{demo/dashboard/static/views.js}; \texttt{demo/dashboard/views.py}; \texttt{demo/ledger/.gitignore}; \texttt{demo/review/.gitignore}; \texttt{demo/review/build\_review.py}; \texttt{demo/review/capture.py}; \texttt{demo/review/demo\_review.tex}; \texttt{demo/review/leave\_behind.tex}; \texttt{tests/dashboard\_views\_smoke.js}; \texttt{tests/test\_dashboard.py}; \texttt{tests/test\_dashboard\_views.py}\par}


\section{Known limitations and pending items}
\begin{longtable}{>{\raggedright\arraybackslash}p{0.30\textwidth}>{\raggedright\arraybackslash}p{0.36\textwidth}>{\raggedright\arraybackslash}p{0.26\textwidth}}
\toprule
\textbf{Item} & \textbf{What resolves it} & \textbf{When}\\\midrule\endhead
The route-at-start decision covers qwen3:14b on evo-x2 and the evo-t2s iGPU only; other device and model pairs show
``not measured'' for it & Route-at-start rows for more models and devices & Not scheduled by this branch\\
evo-t2s local accuracy in the route decision is the R2 session pass rate, not B3 task accuracy & The evo-t2s B5 task
runs & Access window 2026-10-26 to 2026-11-08\\
Hybrid escalation cost against claude-haiku-5-5 is interim (its all-cloud price) & A full-history redo on
claude-haiku-5-5 & Not scheduled by this branch\\
The P1 cost per day and the problem ranking stay ``being measured'' & Their register rows recomputed with the
full-history redo instead of the withheld first redo & When the impact rows are regenerated\\
The fix rate counts 28 inherited-action tasks as not fixed & Stated beside every fix rate; a grader that separates
repairs of earlier local actions would change the count & Not planned for the freeze\\
evo-t2s Intel CPU figures in the trade-off are projected from evo-x2's 4096 tier & The evo-t2s B5 runs & Access
window 2026-10-26 to 2026-11-08 (per row t2s-config-comparison)\\
docs/OCT29\_COST\_\allowbreak SUMMARY.md not present & The file landing on main; the cost panel reads it at request time &
When it lands on main\\
Blade and evo-t2s B5 settings not measured & Runs on those machines; the views pick them up after a cache rebuild &
Next access windows\\
Live escalation is one live fix (n = 1) per demo run: view D sends a single call; the step 4 button runs one
full-history redo turn, which may take several provider requests while it uses tools & By design (USD 10 cap, one call
per run); the full-history redo of all wrong tasks is the measured basis & Not planned\\
Token counts in priced figures use the harness estimate, not Anthropic's tokenizer & Re-pricing on provider usage &
Not planned for the freeze\\
The cache must be rebuilt after register changes & \texttt{py -3.12 -m demo.dashboard.build\_cache} before the demo &
Demo day checklist\\
Backup recording is 3 minutes of a 5-minute talk & Talk over it at the same pace; steps are 36 s each & n/a\\
\bottomrule
\end{longtable}

\appendix
\section{Demo script}
\setlength{\parindent}{0pt}\setlength{\parskip}{4pt}
The script as committed in \texttt{demo/DEMO\_SCRIPT.md}, typeset. Timings: five steps of about one minute each.

The dashboard opens in story mode: five steps, about one minute each, with Next and Back buttons (or the arrow keys)
and a step indicator. After step 5, "Finish: explore freely" (or "Explore freely" at any time) shows the four full
views (A Silent failure, B Device setting, C Cost, D Live escalation) and the older frontier and replay screens.
"Story mode" in the header starts the story again.

Every number on screen is a link: click it to see the register row (docs/NUMBERS\_REGISTER.md) or the result-file line
it came from. Every claim carries a label: measured, priced, not executed, projected, pending flash-attention control,
or not measured. Say the label out loud when you read a number.

The numbers quoted below are the register values at commit time. Read them off the screen on the day: the screen is
rebuilt from the register, this file is not. Steps 1 to 4 run entirely from the offline cache; only step 5 can use
the network, and without it nothing is sent and the cost reads "being measured".

Cost to fix a wrong answer: being measured. The earlier cloud redo (results/cloud\_redo\_v1.jsonl) was a broken
protocol run (thinking consumed the output budget; docs/FINDINGS.md, cloud redo diagnosis), so none of its figures or
answers are shown or quoted anywhere (operator, 2026-10-22). The corrected full-history redo replaces it in the final
pass.

Backup: if the laptop or the browser fails, play demo/review/demo\_story.webm (3 minutes, no audio) and talk over it
with the same script. The two-page leave-behind is demo/review/LEAVE\_BEHIND.pdf. Both still show the earlier redo
figure and are regenerated in the final pass; until then do not use them in front of an audience.

\subsection*{Before the audience arrives (5 minutes, not part of the talk)}

\par\noindent 1. From the repo root on the controller laptop, in the terminal you will present from:

\begin{verbatim}
py -3.12 -m demo.dashboard.build_cache
py -3.12 -m pytest tests/test_dashboard.py tests/test_dashboard_views.py -q
\end{verbatim}

\par\noindent 2. Start the server. Without a live call (step 5 shows the task and the local answer; the cloud cost reads "being
measured"):

\begin{verbatim}
py -3.12 -m demo.dashboard
\end{verbatim}

With the one live call for step 5 (key loaded into this process only, never printed):

\begin{verbatim}
$env:ANTHROPIC_API_KEY = [Environment]::GetEnvironmentVariable('ANTHROPIC_API_KEY','User')
$env:CLOUD_CLIENT_HOST_ROLE = 'controller'; $env:DEMO_CLOUD_LIVE = '1'
py -3.12 -m demo.dashboard
\end{verbatim}

After the demo: Ctrl+C, then \texttt{Remove-\allowbreak{}Item Env:\allowbreak{}\textbackslash{}ANTHROPIC\_\allowbreak{}API\_\allowbreak{}KEY; Remove-\allowbreak{}Item Env:\allowbreak{}\textbackslash{}DEMO\_\allowbreak{}CLOUD\_\allowbreak{}LIVE}.

\par\noindent 3. Open http://127.0.0.1:8765/ : step 1 of 5 shows, no "Error" text. (\texttt{?step=3} opens a given step; \texttt{?explore=1}
skips the story.)

Cloud spend rules for step 5: one provider request per server start (a second click sends nothing), through
src/cloud/client.py with a USD 10 cap (DEMO\_CLOUD\_CAP\_USD may lower it, never raise it) and its own ledger,
demo/ledger/cloud\_ledger\_demo.jsonl (not committed). Any failure (no key, host refused, cap, provider error) shows
"being measured" with the reason on screen. Do not restart the server to get a second live call during the demo.

\subsection*{Step 1 (0:00 to 1:00): the problem}

Click: nothing; step 1 is the opening screen.

Say:
\begin{itemize}\setlength\itemsep{1pt}
\item "Same machine, evo-t2s. Same two models, same 40-turn agent sessions, Ollama defaults, nothing configured. The only
difference is which device Ollama picked."
\item "On the Intel CPU Ollama loads a 4096-token window. Both models silently lost a rule in 5 of 5 sessions, all after
the window filled. On the iGPU it loads 32768: llama3.1:8b lost nothing; qwen3:14b had 1 of 5, and that one also
happens inside the window on its validation run, so 0 after the window filled." Measured (t2s-r2-real-cpu,
t2s-r2-real-igpu).
\item "Calls that surfaced an error to the client: 0 of 800 on each device." Measured (the two result files).
\item "What it does to work: qwen3:14b task results go from 11.3\% wrong at the default window to 42.0\% at 4096, and none of
the 63 wrong results surfaced an error." Measured (x2-evalset-tasks-real-v2-wrong-rate; excess +30.7 points,
session-level 95\% CI +26.0 to +34.7, x2-evalset-tasks-real-v2-ci-units).
\end{itemize}

Do not say: "every model fails at 4096" (step 2 shows why).

\subsection*{Step 2 (1:00 to 2:00): why it is hard}

Click: Next.

Say:
\begin{itemize}\setlength\itemsep{1pt}
\item "The one-line fix, num\_keep, which keeps the system prompt: 30.0\% wrong at 4096. Better than 42.0\%, still far from
11.3\%." Measured (x2-evalset-tasks-real-v2-wrong-rate).
\item "Client-side trimming has a side effect: with trimming, qwen3:14b called the logging tool in 24, 25 and 25 of 40
turns per seed; without it, 40 of 40. Trimming removes the earlier turns that showed the two-call pattern the model
imitates." Measured counts (R2-qwen3-14b-4096-reconciliation); the interpretation is the one in docs/FINDINGS.md.
\item "And it depends on the model: at 4096, llama3.1:8b lost rules in 5 of 5 sessions, qwen3:14b 4 of 5, qwen3:8b 1 of 5;
qwen3:32b 2 of 3; llama3.3:70b and mistral-small3.2:24b resisted rule loss, 0 of 3." Measured (R2-v1b-gated-per-model,
H-scale-verdict, H-scale-family-verdict).
\end{itemize}

\subsection*{Step 3 (2:00 to 3:00): the tool}

Click: Next. The picker opens on evo-x2 / qwen3:14b / multi-turn agent session. Then change Device to "evo-t2s, Intel
CPU". If time allows, "Blade 14".

Say:
\begin{itemize}\setlength\itemsep{1pt}
\item evo-x2 / qwen3:14b: "Window 40960; KV cache q8\_0 with flash attention. Same rule and history scores as the default,
75.4 minutes per session instead of 214.8." The minutes are measured (b5-policy-metrics, cost-cut-config-time).
\item Point at the yellow labels: "The KV and flash attention tiles say pending flash-attention control: we switched both
on together, so we do not yet know which one gives the speedup. When the control run reads, this view recommends
whichever one it shows." Row b5-fa-control is VERIFIED but its run is incomplete ("reading: not read
(incomplete)"). Same for llama3.1:8b (q4\_0). Do not say "q8\_0 makes it faster".
\item evo-t2s CPU / qwen3:14b: "Here nothing we measured holds the rules: the default loses them and trimming did not
recover qwen3:14b, so the tool says not measured instead of inventing a setting." (t2s-r2-real-cpu,
t2s-r2-mitigation-cpu). Switch the model to llama3.1:8b to show the case trimming did recover.
\item Blade: "Not measured. Every field says so."
\end{itemize}

\subsection*{Step 4 (3:00 to 4:00): the money}

Click: Next.

Say:
\begin{itemize}\setlength\itemsep{1pt}
\item "Sending one long agent task whole to the cloud, a 62-step agent task at the median (OpenHands traces): about USD 1.54 on
claude-sonnet-5-5, USD 2.63 at the 90th percentile, assuming a 70\% cache hit rate; about USD 0.08 on
claude-haiku-5-5." Label: priced, not executed (cloud-cost-per-real-task; token counts from traces, prices from the
price file, no API call).
\item Read the line exactly as the screen shows it: "Cost to fix a wrong answer: being measured." Say: "Our first cloud
redo had a protocol fault, so we are re-measuring it; we show no figure until the corrected run is in."
\item "Sending every long task to the cloud is not cheap; escalating only the wrong step needs to know which step went
wrong, and the failure is silent."
\item Do not quote any figure from the earlier redo (cost, seconds, how many it got right) or the P1 cost per day built on
it. Do not quote a router saving: the router was dominated in our own comparison on evo-x2 (b5-verdict: refuted).
\end{itemize}

\subsection*{Step 5 (4:00 to 5:00): one live escalation}

Click: Next, then "Send to the cloud" once.

Say, before clicking:
\begin{itemize}\setlength\itemsep{1pt}
\item "One real wrong answer from qwen3:14b at a 4096 window: book a meeting with the default length from the standing
instructions. The local model set 60 minutes; the instruction gives 50. HTTP 200, no error." Measured
(results/x2\_evalset\_tasks\_real\_v2.jsonl, one click away).
\item Right card: "Cost to fix a wrong answer: being measured." Do not quote the earlier redo of this task.
\end{itemize}

After clicking:
\begin{itemize}\setlength\itemsep{1pt}
\item Live (label measured): "What it cost, how long it took, and what the cloud answered." One call, n = 1, tools offered,
no tool round trips: say "one call", not "this is what fixing costs".
\item Not sent live (label "being measured"): read the reason line ("live calls are off", "no key", or the provider error);
the tiles read "being measured". That is the expected fallback.
\end{itemize}

Close (10 seconds): "A wrong window is silent on every device we measured; the right setting depends on the device and
the model, and we show it only where we measured it. The expensive part is not knowing." Then click "Finish: explore
freely" for questions.

\subsection*{Which claims are measured}

\begin{longtable}{p{0.5\textwidth}p{0.18\textwidth}p{0.24\textwidth}}\toprule
\textbf{Claim on screen} & \textbf{Label} & \textbf{Source} \\ \midrule
CPU loads 4096, iGPU 32768, same models (evo-t2s, Ollama default) & measured & t2s-r2-real-cpu, t2s-r2-real-igpu \\
CPU: 5/5 sessions per model silently lost a rule after the window filled & measured & t2s-r2-real-cpu \\
0 of 800 calls surfaced an error, per device & measured & results/t2s\_r2\_real\_cpu\_v1.jsonl, results/t2s\_r2\_real\_igpu\_v1.jsonl \\
qwen3:14b wrong-task rate 11.3\% default to 42.0\% at 4096, 0 surfaced & measured & x2-evalset-tasks-real-v2-wrong-rate \\
Excess +30.7 points at 4096 (95\% CI +26.0 to +34.7) & measured & x2-evalset-tasks-real-v2-ci-units \\
num\_keep at 4096: 30.0\% wrong & measured & x2-evalset-tasks-real-v2-wrong-rate \\
Trimming: logging tool 24/25/25 of 40 vs 40/40/40 & measured (interpretation from docs/FINDINGS.md) & R2-qwen3-14b-4096-reconciliation \\
Rule loss at 4096 per model, 8B to 70B & measured & R2-v1b-gated-per-model, H-scale-verdict, H-scale-family-verdict \\
evo-x2 qwen3:14b: 40960, 75.4 vs 214.8 min per session & measured & b5-policy-metrics, cost-cut-config-time \\
evo-x2: q8\_0 (qwen3:14b) / q4\_0 (llama3.1:8b) with flash attention as the setting & pending flash-attention control & b5-fa-control (not read yet) \\
evo-t2s CPU qwen3:14b: no setting holds the rules & measured (what failed); setting not measured & t2s-r2-real-cpu, t2s-r2-mitigation-cpu \\
Blade: every setting & not measured & none \\
One long agent task in the cloud: USD 1.54 median, 2.63 p90 (sonnet, 70\% cache hits assumed) & priced, not executed & cloud-cost-per-real-task \\
Cost to fix a wrong answer in the cloud & being measured (no figure; the earlier redo was a broken protocol run) & none until cloud-redo-full-context \\
P1 cost per day for one team, and the problem ranking & being measured (both include the redo cost) & none until re-computed \\
Router savings & not shown (b5-verdict refuted) & b5-verdict \\
Step 5 live answer, USD, seconds & measured, n = 1, this run only & demo/ledger/cloud\_ledger\_demo.jsonl (not committed) \\
Step 5 fallback when not sent live & being measured & none \\
\bottomrule\end{longtable}


\section{Register rows used}
Current values from docs/NUMBERS\_REGISTER.md, verbatim. Not used and not listed: cloud-redo-cost-per-wrong-task (the
first, broken redo run) and impact-v2-P1 and impact-ranking-v2 (built on its cost per wrong task); the demo withholds
them in place.
\paragraph{\texttt{H-scale-family-verdict}} (VERIFIED, n 1 models tested)

{\footnotesize\raggedright overall alternative supported (tested: ['mistral-small3.2:24b']); mistral-small3.2:24b: R4096 0/3 (window exceeded 3/3, gated 0/3), H32768 2/3 (window exceeded 3/3); resists rule loss; history loss with rules holding holds; verdict alternative (resists rule loss)\par}

\paragraph{\texttt{H-scale-verdict}} (VERIFIED, n 2 models tested)

{\footnotesize\raggedright overall mixed (tested: ['qwen3:32b', 'llama3.3:70b']); qwen3:32b: R4096 2/3 (window exceeded 3/3, gated 2/3), H32768 1/3 (window exceeded 3/3); rule loss holds; does not hold; verdict mixed / qwen3-30b-a3b-2507: refused (not tested) (['qwen3-30b-a3b-2507: negative control failed (canary misses 2 at num\_ctx 131072)', 'qwen3-30b-a3b-2507: task tool called correctly on 25/30 baseline turns (\textless{} 90\%, task-tool gate)']) / llama3.3:70b: R4096 0/3 (window exceeded 3/3, gated 3/3), H32768 3/3 (window exceeded 3/3); resists rule loss; history loss with rules holding holds; verdict alternative (resists rule loss)\par}

\paragraph{\texttt{R2-mitigation-v1-verdict}} (VERIFIED, n 12 sessions, 480 turns)

{\footnotesize\raggedright pre-registered verdict at num\_ctx\_4096: refuted; num\_ctx\_8192 llama3.1:8b recovered; num\_ctx\_8192 qwen3:14b recovered; num\_ctx\_4096 llama3.1:8b recovered; num\_ctx\_4096 qwen3:14b not\_recovered (rules below 90\%: rule2\_log\_event\_called)\par}

\paragraph{\texttt{R2-qwen3-14b-4096-reconciliation}} (VERIFIED, n 440 turns)

{\footnotesize\raggedright v1: rules in use 5, Ollama not recorded, harness blob not recorded, call2 off, max tokens 384, client trim None, num\_keep None; per seed 20260901 rule1 40/40, rule2 40/40, rule3 40/40, rule4 20/20, rule5 39/40 / 20260902 rule1 40/40, rule2 40/40, rule3 38/40, rule4 20/20, rule5 40/40 / 20260903 rule1 40/40, rule2 40/40, rule3 37/40, rule4 20/20, rule5 39/40; gated first failure 20260901 5 rule5\_session\_code; 20260902 20 rule3\_no\_zebra; 20260903 5 rule5\_session\_code // v1b: rules in use 5, Ollama 0.34.4, harness blob not recorded, call2 off, max tokens 384, client trim None, num\_keep None; per seed 20260904 rule1 40/40, rule2 40/40, rule3 39/40, rule4 20/20, rule5 40/40 / 20260905 rule1 40/40, rule2 40/40, rule3 40/40, rule4 20/20, rule5 40/40; gated first failure 20260904 30 rule3\_no\_zebra; 20260905 none // mitigation: rules in use 5, Ollama 0.34.4, harness blob not recorded, call2 off, max tokens 384, client trim \{'margin': 0.05\}, num\_keep None; per seed 20260901 rule1 40/40, rule2 24/40, rule3 40/40, rule4 20/20, rule5 39/40 / 20260902 rule1 40/40, rule2 25/40, rule3 40/40, rule4 20/20, rule5 40/40 / 20260903 rule1 40/40, rule2 25/40, rule3 39/40, rule4 20/20, rule5 39/40; gated first failure n/a (not a real-mode run) // num\_keep: rules in use 5, Ollama 0.34.4, harness blob not recorded, call2 off, max tokens 384, client trim None, num\_keep system; per seed 20260901 rule1 40/40, rule2 40/40, rule3 40/40, rule4 20/20, rule5 39/40 / 20260902 rule1 40/40, rule2 40/40, rule3 38/40, rule4 20/20, rule5 40/40 / 20260903 rule1 40/40, rule2 40/40, rule3 37/40, rule4 20/20, rule5 39/40; gated first failure 20260901 5 rule5\_session\_code; 20260902 20 rule3\_no\_zebra; 20260903 5 rule5\_session\_code // num\_keep vs v1 on the same seeds: identical final text 120/120, identical first-call prompt tokens 120/120 // first trimmed rule2 failure seed 20260901 turn 3: call 1 tools trimmed ['lookup\_fact'], untrimmed v1 ['log\_event', 'lookup\_fact']\par}

\paragraph{\texttt{R2-real-v1b-gated-kill}} (VERIFIED, n 75)

{\footnotesize\raggedright gated killed=False; ollama\_default 3/15 silent, 0 after window exceeded; num\_ctx\_32768 1/15 silent, 0 after window exceeded; num\_ctx\_16384 2/15 silent, 2 after window exceeded; num\_ctx\_8192 8/15 silent, 8 after window exceeded; num\_ctx\_4096 10/15 silent, 10 after window exceeded (llama3.1:8b excludes nothing; qwen3:14b excludes tool\_args, recall; qwen3:8b excludes tool\_args, recall)\par}

\paragraph{\texttt{R2-v1b-gated-per-model}} (VERIFIED, n 75)

{\footnotesize\raggedright ollama\_default llama3.1:8b: 3/5 silent, 0 after window exceeded, 3 baseline-type; intact 4/4/4/2/2/2/1/1; ollama\_default qwen3:14b: 0/5 silent, 0 after window exceeded, 0 baseline-type; intact 5/5/5/5/5/0/0/0; ollama\_default qwen3:8b: 0/5 silent, 0 after window exceeded, 0 baseline-type; intact 5/5/5/4/4/0/0/0; num\_ctx\_32768 llama3.1:8b: 1/5 silent, 0 after window exceeded, 1 baseline-type; intact 5/5/5/3/3/0/0/0; num\_ctx\_32768 qwen3:14b: 0/5 silent, 0 after window exceeded, 0 baseline-type; intact 5/5/5/5/0/0/0/0; num\_ctx\_32768 qwen3:8b: 0/5 silent, 0 after window exceeded, 0 baseline-type; intact 5/5/5/4/0/0/0/0; num\_ctx\_16384 llama3.1:8b: 2/5 silent, 2 after window exceeded, 0 baseline-type; intact 5/5/0/0/0/0/0/0; num\_ctx\_16384 qwen3:14b: 0/5 silent, 0 after window exceeded, 0 baseline-type; intact 5/5/0/0/0/0/0/0; num\_ctx\_16384 qwen3:8b: 0/5 silent, 0 after window exceeded, 0 baseline-type; intact 5/5/0/0/0/0/0/0; num\_ctx\_8192 llama3.1:8b: 4/5 silent, 4 after window exceeded, 0 baseline-type; intact 5/0/0/0/0/0/0/0; num\_ctx\_8192 qwen3:14b: 0/5 silent, 0 after window exceeded, 0 baseline-type; intact 5/0/0/0/0/0/0/0; num\_ctx\_8192 qwen3:8b: 4/5 silent, 4 after window exceeded, 0 baseline-type; intact 3/0/0/0/0/0/0/0; num\_ctx\_4096 llama3.1:8b: 5/5 silent, 5 after window exceeded, 0 baseline-type; intact 0/0/0/0/0/0/0/0; num\_ctx\_4096 qwen3:14b: 4/5 silent, 4 after window exceeded, 0 baseline-type; intact 0/0/0/0/0/0/0/0; num\_ctx\_4096 qwen3:8b: 1/5 silent, 1 after window exceeded, 0 baseline-type; intact 0/0/0/0/0/0/0/0\par}

\paragraph{\texttt{T2S-vs-X2-default-ctx}} (VERIFIED, n 3)

{\footnotesize\raggedright \{"qwen3-4b-2507": \{"t2s\_default\_ctx": 4096, "x2\_default\_ctx": 262144\}, "llama3.1:8b": \{"t2s\_default\_ctx": 4096, "x2\_default\_ctx": 131072\}, "qwen3:8b": \{"t2s\_default\_ctx": 4096, "x2\_default\_ctx": 40960\}\}\par}

\paragraph{\texttt{b5-fa-control}} (VERIFIED, n 1 sessions)

{\footnotesize\raggedright llama3.1:8b: (i) f16+FA 44.7 min per session (per seed 20260901: 44.7), (g)/(h) 6.8 min, (a) 43.7 min, window 131072; rule 1.0, hist 0.875; base-env OLLAMA\_FLASH\_ATTENTION ['false']; valid False (no run\_end; sessions completed for seeds [20260901], expected [20260901, 20260902, 20260903]); reading: not read (incomplete) / qwen3:14b: (i) f16+FA None min per session (per seed ), (g)/(h) 75.4 min, (a) 214.8 min, window 40960; rule None, hist None; base-env OLLAMA\_FLASH\_ATTENTION ['false']; valid False (no run\_end; sessions completed for seeds [], expected [20260901, 20260902, 20260903]); reading: not read (incomplete) (tolerance 10\%)\par}

\paragraph{\texttt{b5-kv-recall}} (VERIFIED, n 1 platform files)

{\footnotesize\raggedright evo-x2 llama3.1:8b: f16 ctx 131072 recall 50.0\% (2 calls, truncated 0, errors 1; KV check n/a); q4\_0 ctx 131072 recall 58.3\% (3 calls, truncated 0, errors 0; KV check True); q8\_0 ctx 131072 recall 54.2\% (3 calls, truncated 0, errors 0; KV check True) / evo-x2 qwen3:14b: f16 ctx 40960 recall 100.0\% (3 calls, truncated 0, errors 0; KV check n/a); q4\_0 ctx 40960 recall 100.0\% (3 calls, truncated 0, errors 0; KV check True); q8\_0 ctx 40960 recall 100.0\% (3 calls, truncated 0, errors 0; KV check True) / evo-t2s: not run yet / blade: not run yet\par}

\paragraph{\texttt{b5-policy-metrics}} (VERIFIED, n 1 platform files)

{\footnotesize\raggedright evo-x2 llama3.1:8b: a rule 100.0\% hist 95.8\% 43.7 min \$0.0000; b rule 100.0\% hist 95.8\% 43.7 min \$0.0000 (=a); c rule 100.0\% hist 95.8\% 43.6 min \$0.0000; d rule 100.0\% hist 95.8\% 43.5 min \$0.0000; e rule 100.0\% hist 95.8\% 43.7 min \$0.0000 (=a); f rule 100.0\% hist 95.8\% 42.6 min \$0.0000; g rule 100.0\% hist 95.8\% 6.9 min \$0.0000; h rule 100.0\% hist 100.0\% 6.6 min \$0.0000; frontier ['h']; complete True / evo-x2 qwen3:14b: a rule 100.0\% hist 62.5\% 214.8 min \$0.0000; b rule 100.0\% hist 62.5\% 214.8 min \$0.0000 (=a); c rule 100.0\% hist 62.5\% 178.7 min \$0.0000; d rule 100.0\% hist 62.5\% 135.1 min \$0.8419; e rule 100.0\% hist 62.5\% 214.8 min \$0.0000 (=a); f rule 100.0\% hist 62.5\% 207.8 min \$0.0000; g rule 100.0\% hist 62.5\% 75.4 min \$0.0000; h rule 100.0\% hist 62.5\% 75.4 min \$0.0000; frontier ['g']; complete True / evo-t2s: not run yet / blade: not run yet\par}

\paragraph{\texttt{b5-verdict}} (VERIFIED, n 3 platforms)

{\footnotesize\raggedright refuted (router dominated in evo-x2/llama3.1:8b, evo-x2/qwen3:14b); evo-x2 complete (router on frontier in every complete cell: False, fixed on frontier in every cell: []); evo-t2s not run (router on frontier in every complete cell: None, fixed on frontier in every cell: None); blade not run (router on frontier in every complete cell: None, fixed on frontier in every cell: None)\par}

\paragraph{\texttt{cloud-cost-per-real-task}} (VERIFIED, n 636 trajectories)

{\footnotesize\raggedright USD per real agent task, median / p90, priced, not executed (no API call): nebius-rebench-openhands (n=36 tasks, median 62 steps; billed input tokens median 1,721,728 p90 3,021,217; billed output median 3,348 p90 4,806): claude-haiku-5-5: none \$0.1738 / \$0.3046, realistic \$0.0826 / \$0.1421, best \$0.0255 / \$0.0401; claude-sonnet-5-5: none \$3.4761 / \$6.0913, realistic \$1.5367 / \$2.6348, best \$0.3378 / \$0.5060 \textbar{} SWE-Gym (n=600 tasks, median 6 steps; billed input tokens median 9,566 p90 422,654; billed output median 324 p90 1,620): claude-haiku-5-5: none \$0.0012 / \$0.0430, realistic \$0.0008 / \$0.0216, best \$0.0005 / \$0.0079; claude-sonnet-5-5: none \$0.0234 / \$0.8610, realistic \$0.0152 / \$0.4040, best \$0.0094 / \$0.1182; cases: none = all input uncached, realistic = ASSUMPTION cache hit rate 0.7 of the re-sent prefix (rest written at the 5-minute cache-write premium), best = every re-sent prefix read from cache, only new tokens written; each step re-sends its full prompt; output tokens per step are not in the trace data, STAND-IN 54 tokens per step (median per-call reply, results/x2\_r2\_real\_v1.jsonl, n=1440 calls); prices results/cloud\_prices\_anthropic\_20261014.json Standard; token basis qwen tokenizer, not Anthropic's (assumption)\par}

\paragraph{\texttt{cloud-redo-cost-per-wrong-task}} (VERIFIED, n 197 wrong tasks)

{\footnotesize\raggedright mode executed; model claude-sonnet-5-5 (results/cloud\_prices\_anthropic\_20261014.json), prompt caching on (cold cache per task); USD per wrong task mean \$0.0201, median \$0.0122, total \$3.9520 (all input uncached: \$13.8989); per tier: ollama\_default n=17 mean \$0.0364 median \$0.0297 total \$0.6193; ollama\_ctx\_16384 n=19 mean \$0.0373 median \$0.0356 total \$0.7084; ollama\_ctx\_8192 n=53 mean \$0.0151 median \$0.0112 total \$0.8007; ollama\_ctx\_4096 n=63 mean \$0.0149 median \$0.0117 total \$0.9356; ollama\_ctx\_4096\_numkeep n=45 mean \$0.0197 median \$0.0102 total \$0.8880; local regrade on the rebuilt environment reproduces end\_state\_ok false for 197/197 (100.0\%); token basis is the harness estimate, not Anthropic's tokenizer; [cloud fix rate withheld in the demo: under diagnosis], mean 6.89 s per task\par}

\paragraph{\texttt{cost-cut-config-time}} (VERIFIED, n 2 models)

{\footnotesize\raggedright MEASURED; evo-x2 llama3.1:8b: default (a) 43.7 min per session, rule 100.0\% hist 95.8\%; fastest measured setting (h) kv\_q4\_0 6.6 min, rule 100.0\% hist 100.0\%, window 131072; saved 37.1 min per session (85\%), over 3 seeds / evo-x2 qwen3:14b: default (a) 214.8 min per session, rule 100.0\% hist 62.5\%; fastest measured setting (g) kv\_q8\_0 75.4 min, rule 100.0\% hist 62.5\%, window 40960; saved 139.4 min per session (65\%), over 3 seeds; flash-attention control included\par}

\paragraph{\texttt{impact-P3}} (VERIFIED, n 3)

{\footnotesize\raggedright claude-sonnet-5-5 standard caching on, sessions shaped like llama3.1:8b's: status measured; measured USD/day 3.60, assumed USD/day 0.00, range USD/day 3.60 to 71.96; per session: re-sent history 1,798,884 tok, USD 0.1799 caching on, 3.5978 no caching; whole session to the cloud USD 0.3358 caching on, 3.7302 no caching (n=5); assumptions: team size = 10 (assumption (profile JSON)); agent sessions per developer per day = 4 (assumption (profile JSON)); cloud session share = 0.5 (assumption (profile JSON)); cloud model = claude-sonnet-5-5 standard, caching on (results/cloud\_prices\_anthropic\_20261014.json (official Anthropic pricing page)); token basis = o200k\_base counts (assumption: Anthropic tokenizer counts not measured); re-sent history is a cacheable prefix = R2 sessions are append-only (best case for caching; real histories not measured); confidence high\par}

\paragraph{\texttt{impact-hour-value}} (VERIFIED, n 1)

{\footnotesize\raggedright USD 64.44/h (\$64.44 per hour; Quick Facts: Software Developers, Quality Assurance Analysts, and Testers, row "2025 Median Pay"; May 2025; https://www.bls.gov/ooh/computer-and-information-technology/software-developers.htm; fetched 2026-10-14)\par}

\paragraph{\texttt{impact-ranking-v2}} (VERIFIED, n 7 rows)

{\footnotesize\raggedright one profile: evo-x2 qwen3:14b ollama\_ctx\_4096, team 10 x 4 sessions/day, cloud claude-sonnet-5-5 caching on (ASSUMED volumes; hour value cited); ranking by main-case USD/day (P2: machine capacity; the other rows: time at the full hour value; a PENDING row is not ranked): P1 \textgreater{} P5 \textgreater{} P2 \textgreater{} P3 \textgreater{} P4; unranked: P6, P7; rows: P1 132.31 (range 5.47 to 132.31; mixed, cloud part measured); P5 11.04 (range 0.00 to 11.04; mixed); P2 5.48 (range 5.48 to 1,503.10; mixed; machine capacity: qwen3:14b default (a) 214.8 min per session = 6.70 sessions/machine/day vs (g) kv\_q8\_0 75.4 min = 19.11 (MEASURED minutes, cost-cut-config-time; 24 h / session minutes, one session at a time, no measured concurrency row); machines for the profile's sessions/day (ASSUMED volume): 6 default vs 3 best (ceil), 3 saved x USD 2,000 (impact-ref-machine-price, ASSUMPTION, must be verified) / 3-year life (ASSUMED) = USD 5.48/day; upper bound: waiting time 139.4 min saved per session x hour value x interactive-session share 0.25 (ASSUMED) = USD 1,497.62/day on top; the flash-attention control (row b5-fa-control) is in the comparison; its reading says whether the saving is flash attention or KV precision); P3 3.70 (range 3.70 to 74.10; measured); P4 0.00 (range 0.00; mixed); P6 not measured (range n/a; not measured); P7 not measured (range n/a; not measured); split: measured P3, mixed P1/P2/P4/P5, not measured P6/P7; PENDING inputs: none\par}

\paragraph{\texttt{impact-total}} (VERIFIED, n 7)

{\footnotesize\raggedright status quo ollama\_default, claude-sonnet-5-5 caching on: measured-only USD/day 3.60 (USD/month 75.55); range USD/day 3.60 to 192.81 (USD/month 75.55 to 4,049.08; low = measured only, nobody waits for compute; high = compute time at the full hour value and cloud without caching) over P2, P3, P4, P5, P6, P7; not added (assumed or not measured): P1; ranking (caching on, time at its upper bound) P5 \textgreater{} P6 \textgreater{} P2 \textgreater{} P3 \textgreater{} P4 \textgreater{} P7 \textgreater{} P1; per row status P1 assumed, P2 mixed, P3 measured, P4 mixed, P5 mixed, P6 mixed, P7 measured\par}

\paragraph{\texttt{impact-v2-P1}} (VERIFIED, n 4 tiers)

{\footnotesize\raggedright profile evo-x2 qwen3:14b ollama\_ctx\_4096, 40 sessions/day (ASSUMED volume), hour value USD 64.44 (impact-hour-value); status mixed, cloud part measured; profile cell USD/day 5.47 to 132.31; ollama\_ctx\_16384: excess +1.3 points (session-level 95\% CI -1.3 to +4.0), 30 task turns/session, 16.0 excess wrong tasks/day, wrong-turn wall 12.5 s; time USD/day 0.00 to 3.57 (CI 0.00 to 10.72); cloud redo USD/day 0.60 (USD 0.0373 per wrong task, measured; CI 0.00 to 1.79); ollama\_ctx\_8192: excess +24.0 points (session-level 95\% CI +18.7 to +29.3), 30 task turns/session, 288.0 excess wrong tasks/day, wrong-turn wall 30.5 s; time USD/day 0.00 to 157.24 (CI 122.30 to 192.19); cloud redo USD/day 4.35 (USD 0.0151 per wrong task, measured; CI 3.38 to 5.32); ollama\_ctx\_4096: excess +30.7 points (session-level 95\% CI +26.0 to +34.7), 30 task turns/session, 368.0 excess wrong tasks/day, wrong-turn wall 19.3 s; time USD/day 0.00 to 126.85 (CI 107.55 to 143.39); cloud redo USD/day 5.47 (USD 0.0149 per wrong task, measured; CI 4.63 to 6.18); ollama\_ctx\_4096\_numkeep: excess +18.7 points (session-level 95\% CI +16.0 to +21.3), 30 task turns/session, 224.0 excess wrong tasks/day, wrong-turn wall 12.2 s; time USD/day 0.00 to 48.84 (CI 41.87 to 55.82); cloud redo USD/day 4.42 (USD 0.0197 per wrong task, measured; CI 3.79 to 5.05); measured: excess rate and its CI (x2-evalset-tasks-real-v2-wrong-rate, x2-evalset-tasks-real-v2-ci-units), task turns per session, wrong-turn wall time; measured: cloud redo USD per wrong task, per tier (cloud-redo-cost-per-wrong-task); assumed: volume, who waits (0 to full hour value), a wrong task costs one cloud redo plus the wasted turn (time to notice not counted)\par}

\paragraph{\texttt{num-keep-guardrail-test}} (VERIFIED, n 120 turns)

{\footnotesize\raggedright llama3.1:8b num\_ctx 4096 num\_keep=system: S per seed \{20260901: 601, 20260902: 594, 20260903: 597\}; observed context shifts 105 with n\_keep [595, 598, 602], task n\_keep [594, 597, 601], token cuts 7 with keep [594, 601], --keep [], Ollama message-drop calls 224 of 239; rules in use: rule1 94.2\% (n=120, first failure 31/-/5); rule3 100.0\% (n=120, first failure -/-/-); rule4 100.0\% (n=60, first failure -/-/-); canary misses C 6, H 24 of 24 checks; error turns 2; verdict: one-line fix\par}

\paragraph{\texttt{num-keep-qwen3-14b}} (VERIFIED, n 120 turns)

{\footnotesize\raggedright qwen3:14b num\_ctx 4096 num\_keep=system: S per seed \{20260901: 752, 20260902: 747, 20260903: 750\}; observed context shifts 0 with n\_keep [], task n\_keep [747, 750, 752], token cuts 0 with keep [], --keep [], Ollama message-drop calls 228 of 240; rules in use: rule1 100.0\% (n=120, first failure -/-/-); rule2 100.0\% (n=120, first failure -/-/-); rule3 95.8\% (n=120, first failure -/20/20); rule4 100.0\% (n=60, first failure -/-/-); rule5 98.3\% (n=120, first failure 5/-/5); canary misses C 1, H 24 of 24 checks; error turns 0; verdict: inconclusive (flag effect not observed)\par}

\paragraph{\texttt{t2s-r2-mitigation-cpu}} (VERIFIED, n 6 sessions, 240 turns)

{\footnotesize\raggedright overall refuted; num\_ctx\_4096 llama3.1:8b recovered; num\_ctx\_4096 qwen3:14b not\_recovered (rules below 90\%: rule2\_log\_event\_called)\par}

\paragraph{\texttt{t2s-r2-real-cpu}} (VERIFIED, n 10 sessions)

{\footnotesize\raggedright llama3.1:8b: loaded window [4096], first over window turn [3]; gated 5/5 silent, 5 after window exceeded, 0 baseline-type (gate excludes recall); intact 0/0/0/0/0/0/0/0 / qwen3:14b: loaded window [4096], first over window turn [3]; gated 5/5 silent, 5 after window exceeded, 0 baseline-type (gate excludes recall); intact 0/0/0/0/0/0/0/0; pre-registered kill criterion killed=False, ollama\_default 10/10 silent\par}

\paragraph{\texttt{t2s-r2-real-igpu}} (VERIFIED, n 10 sessions)

{\footnotesize\raggedright llama3.1:8b: loaded window [32768], first over window turn [23]; gated 0/5 silent, 0 after window exceeded, 0 baseline-type (gate excludes nothing); intact 5/5/5/4/4/0/0/0 / qwen3:14b: loaded window [32768], first over window turn [20, 21]; gated 1/5 silent, 0 after window exceeded, 1 baseline-type (gate excludes recall); intact 4/4/4/4/0/0/0/0; pre-registered kill criterion killed=False, ollama\_default 5/10 silent\par}

\paragraph{\texttt{x2-evalset-tasks-real-v2-ci-units}} (VERIFIED, n 25 sessions, 750 task turns)

{\footnotesize\raggedright qwen3:14b wrong-task excess over default; ollama\_ctx\_16384: excess +1.3 points; session-level 95\% CI -1.3 to +4.0 (5 vs 5 sessions); turn-level 95\% CI -6.0 to +8.7 (150 vs 150 task turns); ollama\_ctx\_8192: excess +24.0 points; session-level 95\% CI +18.7 to +29.3 (5 vs 5 sessions); turn-level 95\% CI +14.7 to +33.3 (150 vs 150 task turns); ollama\_ctx\_4096: excess +30.7 points; session-level 95\% CI +26.0 to +34.7 (5 vs 5 sessions); turn-level 95\% CI +21.3 to +40.0 (150 vs 150 task turns); ollama\_ctx\_4096\_numkeep: excess +18.7 points; session-level 95\% CI +16.0 to +21.3 (5 vs 5 sessions); turn-level 95\% CI +10.0 to +27.3 (150 vs 150 task turns) (10000 resamples each, seed 20261014; the paper uses the session-level CI)\par}

\paragraph{\texttt{x2-evalset-tasks-real-v2-mistral-small32-wrong-rate}} (PENDING, n n/a)

{\footnotesize\raggedright not computed: data file(s) not synced yet: results/x2\_evalset\_tasks\_real\_v2\_mistral\_small32.jsonl\par}

\paragraph{\texttt{x2-evalset-tasks-real-v2-qwen3-8b-wrong-rate}} (VERIFIED, n 0 task turns)

{\footnotesize\raggedright run finished False; qwen3:8b: refused (qwen3:8b: task tool called correctly on 61/81 baseline turns (\textless{} 90\%, task-tool gate); qwen3:8b: task end-state gate failed (65/90 \textless{} 90\%); qwen3:8b: strict preflight, gates not passed: ['baseline', 'task\_tool', 'task\_end\_state']) (bootstrap 10000 resamples, unit session, stratified by tier, seed 20261014)\par}

\paragraph{\texttt{x2-evalset-tasks-real-v2-wrong-rate}} (VERIFIED, n 750 task turns)

{\footnotesize\raggedright run finished True; llama3.1:8b: refused (llama3.1:8b: negative control failed (canary misses 26 at num\_ctx 131072); llama3.1:8b: not in this plan's model list ['qwen3:14b']; llama3.1:8b: task end-state gate failed (72/90 \textless{} 90\%)) / qwen3:14b: valid True; ollama\_default: wrong 17/150 (11.3\%) (silent 17, surfaced 0), loaded context [40960]; ollama\_ctx\_16384: wrong 19/150 (12.7\%) (silent 19, surfaced 0), loaded context [16384]; ollama\_ctx\_8192: wrong 53/150 (35.3\%) (silent 53, surfaced 0), loaded context [8192]; ollama\_ctx\_4096: wrong 63/150 (42.0\%) (silent 63, surfaced 0), loaded context [4096]; ollama\_ctx\_4096\_numkeep: wrong 45/150 (30.0\%) (silent 45, surfaced 0), loaded context [4096], num\_keep per seed \{20260901: 3369, 20260902: 3369, 20260903: 3364, 20260904: 3369, 20260905: 3367\}; ollama\_ctx\_16384 excess over default +1.3 points (95\% CI -1.3 to +4.0, 5 vs 5 sessions); reading: no excess shown; ollama\_ctx\_8192 excess over default +24.0 points (95\% CI +18.7 to +29.3, 5 vs 5 sessions); reading: ollama\_ctx\_8192 raises the wrong-task rate; ollama\_ctx\_4096 excess over default +30.7 points (95\% CI +26.0 to +34.7, 5 vs 5 sessions); reading: ollama\_ctx\_4096 raises the wrong-task rate; ollama\_ctx\_4096\_numkeep excess over default +18.7 points (95\% CI +16.0 to +21.3, 5 vs 5 sessions); reading: ollama\_ctx\_4096\_numkeep raises the wrong-task rate (bootstrap 10000 resamples, unit session, stratified by tier, seed 20261014)\par}

\paragraph{\texttt{x2-scale-real-default-tier}} (VERIFIED, n 3 sessions)

{\footnotesize\raggedright llama3.3:70b: no default-tier session / qwen3:32b: 3 default-tier sessions completed (seed 20260901: 40 turns, 6.18 h, loaded [40960], first over window 24; seed 20260902: 40 turns, 6.01 h, loaded [40960], first over window 25; seed 20260903: 40 turns, 4.66 h, loaded [40960], first over window 26)\par}

\paragraph{\texttt{x2-v3-usability}} (VERIFIED, n 121)

{\footnotesize\raggedright per-call timeout 900 s (harness DEFAULT\_CALL\_TIMEOUT\_S); timeout\_latency/done (share) [context tiers]: llama3.1:8b/llama\_server 40/340 (11.8\%) [longdoc\_120000 x10, longdoc\_64000 x10, longdoc\_80000 x10, longdoc\_96000 x10]; llama3.1:8b/ollama\_default 10/340 (2.9\%) [longdoc\_120000 x10]; qwen3-14b/llama\_server 0/340 (0.0\%); qwen3-14b/ollama\_default 0/340 (0.0\%); qwen3-30b-a3b/llama\_server 10/340 (2.9\%) [longdoc\_120000 x10]; qwen3-30b-a3b/ollama\_default 0/340 (0.0\%); qwen3-32b/llama\_server 11/100 (11.0\%) [longdoc\_32000 x10, trace\_mix\_24000 x1]; qwen3-32b/ollama\_default 0/100 (0.0\%); qwen3-4b-2507/llama\_server 40/340 (11.8\%) [longdoc\_120000 x10, longdoc\_64000 x10, longdoc\_80000 x10, longdoc\_96000 x10]; qwen3-4b-2507/ollama\_default 10/340 (2.9\%) [longdoc\_120000 x10]; qwen3-8b/llama\_server 0/340 (0.0\%); qwen3-8b/ollama\_default 0/340 (0.0\%)\par}


\end{document}
